\documentclass[10pt,a4paper]{amsart}
\usepackage[margin=19mm]{geometry}
\usepackage{amsmath,amssymb,mathtools}
\usepackage[scr=rsfs,cal=euler]{mathalfa}
\usepackage{xfrac}
\usepackage[unicode,hidelinks,bookmarks=false]{hyperref}
\newcommand{\ie}{i.e.\ }
\newcommand{\cf}{cf.\ }
\newcommand{\resp}{respectively\ }
\newcommand{\Subsection}{\subsection}
\providecommand{\nobreakdash}{\nobreak\hbox{-}}
\setlength{\emergencystretch}{2em}
\multlinegap0pt
\usepackage{bm}
\usepackage{bbm}
\usepackage{dsfont}
\usepackage{xspace}
\usepackage{cancel}
\usepackage{mathtools}
\usepackage{amsthm}
\makeatletter
\@ifundefined{theorem}{\newtheorem{theorem}{Theorem}}{}
\@ifundefined{lemma}{\newtheorem{lemma}[theorem]{Lemma}}{}
\makeatother
\newcommand{\R}{\mathbb{R}}
\title[Boundary Value Problems in graph Lipschitz domains in the pl]{Boundary Value Problems in graph Lipschitz domains in the plane with $A\_{\infty}$-measures on the boundary}
\author{Fernando Ballesta-Yag\"ue, Mar\'ia J. Carro}
\date{Source: arXiv:2506.23961v2 (CC BY 4.0)}
\begin{document}
\maketitle

\noindent\emph{Source and licence.} Boundary Value Problems in graph Lipschitz domains in the plane with $A_{\infty}$-measures on the boundary. Version arXiv:2506.23961v2 (CC BY 4.0), distributed under the Creative Commons Attribution 4.0 International licence, as recorded in the arXiv licence metadata for this exact version and shown on the abstract page https://arxiv.org/abs/2506.23961v2. Full rights notice: licenses/arxiv-2506.23961-CC-BY-4.0.txt.\par\smallskip

\noindent\textit{Editorial scope.} Counted unit: the Lemma whose source label is \texttt{lemma:lemma1RangeSolvability} in the article (source0.tex lines 1458--1460, source label \texttt{lemma:lemma1RangeSolvability}), delivered together with the article's own complete original proof of it (source0.tex lines 1463--1521, one \texttt{proof} environment of 59 lines). No mathematics is written, repaired or re-derived: statement and proof are the article's own text line for line. Required context retained uncounted: lemma:improvingA1 (lines 497--509). Every definition, display and label of those lines is retained unchanged. Omitted from this package and not redistributed: 1--496, 510--1457, 1461--1462, 1522--3209. Nothing inside a retained excerpt is omitted or altered: every retained body is delivered exactly as the article's lines are written, so no omission receipt of any kind accompanies this package. Citations: the retained excerpts carry no citation command at all, so no bibliography entry of the article is transcribed and no reference is redistributed. Reference adaptation: none was needed and none was made; every reference inside the delivered unit resolves inside it, so no unresolved reference or citation is produced. Printed slips kept literal: no printed word, symbol, hypothesis or number of the counted statement or of its proof was altered. Layout and class: the article's own class is \texttt{amsart} with packages including graphicx, amsmath, amssymb, amsfonts, mathtools, amsthm, dsfont, hyperref, mathrsfs, anysize, tikz, bm, enumitem; the wrapper is the lane's pinned \texttt{amsart} editorial wrapper at 10pt on A4, which restates the theorem-like environments the retained text needs and adds the article's own macro definitions that the retained text needs (\texttt{\textbackslash R}), with extra wrapper packages (bm, bbm, dsfont, xspace, cancel, mathtools). The delivered numbering is the unit's own. Assets: no retained span contains a figure environment, a graphics-inclusion command, a PDF-page inclusion, a table or a TikZ picture; no image file is copied into this package, the package declares no asset dependency and no image file is redistributed with it. Licence: version arXiv:2506.23961v2 of this article is distributed under the Creative Commons Attribution 4.0 International licence, as recorded in the arXiv licence metadata for this exact version and shown on the abstract page https://arxiv.org/abs/2506.23961v2. A full rights notice is delivered at licenses/arxiv-2506.23961-CC-BY-4.0.txt. Characters: every retained excerpt is pure ASCII, so the delivered engine, XeLaTeX with its default Latin Modern text font, reports no missing character.
\begin{lemma}\label{lemma:improvingA1}    Let $p\in [1,\infty)$ and $u\in A_p$. 
\begin{enumerate}
\item  Then there exists $r>1$ small enough such that $u^r\in A_p$.
\item There exist $A_1$ weights $w_1$ and $w_2$ such that
\[
w=w_1 w_2^{1-p} .
\]
\end{enumerate}


Moreover, if $w_0,w_1\in A_p$ and $\alpha_0,\alpha_1>0$ satisfy $\alpha_0+\alpha_1\leq 1$, then $w_0^{\alpha_0}w_1^{\alpha_1}\in A_p$.

\end{lemma}

\begin{lemma}\label{lemma:lemma1RangeSolvability}
    Let $\nu\in A_{\infty}(\Lambda)$. If $R(\nu)\neq\varnothing$, then there are at least three points in $R(\nu)$. More specifically, given $p\in R(\nu)$, there exist $q_{-},q_{+}\in R(\nu)$ with $q_{-}<p<q_{+}$. 
\end{lemma}

\begin{proof}
    Let $p\in R(\nu)$. This means that $|\Phi'|^{-p}\Phi(\nu)\in A_p$, which can be rewritten as
    \[
    \nu(\Phi(x))|\Phi'(x)|^{1-p}\in A_p(\R).
    \]
    Therefore, by the Factorization Theorem,  
    \begin{equation}\label{eq:fact1_rangeSolvability}
    \exists u_0,u_1\in A_1(\R): \nu(\Phi(x))|\Phi'(x)|^{1-p}=u_0 u_1^{1-p}.    
    \end{equation}
    We want to check that there exists $q\neq p$ such that $q\in R(\nu)$, i.e., such that $\nu(\Phi(x))|\Phi'(x)|^{1-q}\in A_q$. Again, by the Factorization Theorem, this is equivalent to the existence of $v_0,v_1\in A_1(\R)$ such that
    \begin{equation}\label{eq:fact2_rangeSolvability}
        \nu(\Phi(x))|\Phi'(x)|^{1-p}=v_0 v_1^{1-q}.
    \end{equation}


    We try to use \eqref{eq:fact1_rangeSolvability} to obtain \eqref{eq:fact2_rangeSolvability}. We can write
    \begin{equation}\label{eq:ident1_rangeSolvability}
    \nu(\Phi(x))|\Phi'(x)|^{1-q}
    =
    \nu(\Phi(x))|\Phi'(x)|^{1-p}|\Phi'(x)|^{p-q}
    =
    u_0 u_1^{1-p}|\Phi'(x)|^{p-q}.    
    \end{equation}
    Since $|\Phi'(x)|\in A_2(\R)$, we can use once more the Factorization Theorem to obtain that there exists $\widehat{u}_0,\widehat{u}_1\in A_1(\R)$ such that
    \[
    |\Phi'(x)|
    =
    \widehat{u}_0\widehat{u}_1^{-1}.
    \]
    Substituting this into \eqref{eq:ident1_rangeSolvability}, we have
    \begin{equation}\label{eq:ident2_rangeSolvability}
    \nu(\Phi(x))|\Phi'(x)|^{1-q}
    =
    u_0 u_1^{1-p}(\widehat{u}_0\widehat{u}_1^{-1})^{p-q}
    =
    u_0 u_{1}^{1-p} \widehat{u}_{0}^{p-q} \widehat{u}_{1}^{q-p}.    
    \end{equation}

We have to distinguish the two cases $p-q>0$ and $p-q<0$. We just give the details of the case $p>q$, because the other one can be settled in the same way. So assume $p>q$. Then
    \begin{equation}\label{eq:case1_rangeSolvability}
        u_0 u_1^{1-p} \widehat{u}_{0}^{p-q}\widehat{u}_{1}^{q-p}
        =
        (u_0\widehat{u}_0^{p-q})(u_1^{1-p}\widehat{u}_1^{q-p})
        =
        (u_0\widehat{u}_0^{p-q})(u_1^{\frac{1-p}{1-q}}\widehat{u}_{1}^{\frac{q-p}{1-q}})^{1-q},
    \end{equation}
    so it is enough to prove that $u_0\widehat{u}_0^{p-q}\in A_1(\R)$ and $u_1^{\frac{1-p}{1-q}}\widehat{u}_{1}^{\frac{q-p}{1-q}}\in A_1(\R)$ for some value of $q<p$.

    We have that $u_0\widehat{u}_0^{p-q}\in A_1(\R)$ if  $p-q$ is small enough. Indeed, by Lemma \ref{lemma:improvingA1}, $u_0^{r}\in A_1(\R)$ for some $r>1$ small enough. So we can write
    \[
    u_0\widehat{u}_0^{p-q}
    =
    (u_0^{r})^{\frac{1}{r}}\widehat{u}_{0}^{p-q},
    \]
    which is in $A_1(\R)$ if $\frac{1}{r}+p-q\leq 1$ by Lemma \ref{lemma:improvingA1}. 


    We also have that $u_1^{\frac{1-p}{1-q}}\widehat{u}_{1}^{\frac{q-p}{1-q}}\in A_1(\R)$ if $p-q$ is small enough. The reasoning is analogous using that  there exists $r>1$ sufficiently small such that  $u_1^r\in A_1(\R)$. 
\end{proof}

\end{document}
